% GENERATED FILE. Edit canonical lessons, then run npm run generate:books.
% canonical-source-hash: fnv1a64:2d1a7cf429fb3443
% canonical-lessons: MW-C185-dhaniyo, MW-C185-rai, MW-C185-methi, MW-C185-hig, MW-C185-saumph

\chapter{Spices}
\label{ch:mw-spices}
\hlchaptermodality{185}

\begin{chapteropening}
\textbf{By the end of this chapter:} \emph{I can name the spices in my kitchen.}

Complete the last lesson of chapter 185: I can name the spices in my kitchen.
\end{chapteropening}

\section[\texorpdfstring{\mw{धणियो} (dhaṇiyo)}{धणियो (dhaṇiyo)}]{\mw{धणियो} (dhaṇiyo) — coriander}

\label{lesson:MW-C185-dhaniyo}



\begin{quote}
Before the new one: say the Marwadi for copper, then the Marwadi for iron.
\end{quote}

\subsection*{You'll want to know: \mw{धणियो}}
\textbf{\mw{धणियो}} — \emph{dhaṇiyo} — \textquotedblleft{}coriander\textquotedblright{}.

\begin{quote}
\mw{सबजी} \mw{मांय} \mw{धणियो} \mw{घालो।}
\end{quote}

\begin{quote}
\emph{sabjī mā̃y dhaṇiyo ghālo.}
\end{quote}

\begin{quote}
Put coriander in the vegetables.
\end{quote}

A masculine noun.

\begin{cousinweb}
\textbf{\mw{धणियो}} (\emph{dhaṇiyo}) is inherited from Sanskrit \emph{dhānaka}, \textquotedblleft{}coriander\textquotedblright{}; CDIAL lists Hindi \emph{dhaniyā} (CDIAL 6776). CDIAL gives no Marwari form in this entry.
\end{cousinweb}

\begin{grammarlens}[title={What you've built}]
One more word for this chapter.
\end{grammarlens}

\subsection*{Your turn}
\begin{itemize}
\raggedright
  \item \emph{Say it:} \emph{dhaṇiyo}
  \item \emph{Say it:} \emph{dhaṇiyo}, once more
  \item \emph{Say it:} \emph{sabjī mā̃y dhaṇiyo ghālo.} — Put coriander in the vegetables.
  \item \emph{Recall:} say the Marwadi for copper, then the Marwadi for iron, then say \emph{dhaṇiyo} again
\end{itemize}

\begin{tcolorbox}[breakable,colback=teal!4,colframe=teal!35!black,title={Before you move on}]
What does \mw{धणियो} mean? (\textquotedblleft{}Coriander\textquotedblright{}.) Say it once more.
\end{tcolorbox}
\section[\texorpdfstring{\mw{राई} (rāī)}{राई (rāī)}]{\mw{राई} (rāī) — mustard seed}

\label{lesson:MW-C185-rai}



\begin{quote}
Before the new one: say the Marwadi for iron, then the Marwadi for coriander.
\end{quote}

\subsection*{You'll want to know: \mw{राई}}
\textbf{\mw{राई}} — \emph{rāī} — \textquotedblleft{}mustard seed\textquotedblright{}.

\begin{quote}
\mw{मां} \mw{तेल} \mw{मांय} \mw{राई} \mw{घालै} \mw{है।}
\end{quote}

\begin{quote}
\emph{mā̃ tel mā̃y rāī ghālai hai.}
\end{quote}

\begin{quote}
Mother puts mustard seed in the oil.
\end{quote}

A feminine noun.

\begin{cousinweb}
\textbf{\mw{राई}} (\emph{rāī}) is inherited from Sanskrit \emph{rājikā}, \textquotedblleft{}mustard\textquotedblright{}; CDIAL lists Hindi and Gujarati \emph{rāī} (CDIAL 10688). CDIAL gives no Marwari form in this entry.
\end{cousinweb}

\begin{grammarlens}[title={What you've built}]
One more word for this chapter.
\end{grammarlens}

\subsection*{Your turn}
\begin{itemize}
\raggedright
  \item \emph{Say it:} \emph{rāī}
  \item \emph{Say it:} \emph{rāī}, once more
  \item \emph{Say it:} \emph{mā̃ tel mā̃y rāī ghālai hai.} — Mother puts mustard seed in the oil.
  \item \emph{Recall:} say the Marwadi for iron, then the Marwadi for coriander, then say \emph{rāī} again
\end{itemize}

\begin{tcolorbox}[breakable,colback=teal!4,colframe=teal!35!black,title={Before you move on}]
What does \mw{राई} mean? (\textquotedblleft{}Mustard seed\textquotedblright{}.) Say it once more.
\end{tcolorbox}
\section[\texorpdfstring{\mw{मेथी} (methī)}{मेथी (methī)}]{\mw{मेथी} (methī) — fenugreek}

\label{lesson:MW-C185-methi}



\begin{quote}
Before the new one: say the Marwadi for coriander, then the Marwadi for mustard seed.
\end{quote}

\subsection*{You'll want to know: \mw{मेथी}}
\textbf{\mw{मेथी}} — \emph{methī} — \textquotedblleft{}fenugreek\textquotedblright{}.

\begin{quote}
\mw{मेथी} \mw{कड़वी} \mw{है।}
\end{quote}

\begin{quote}
\emph{methī kaṛvī hai.}
\end{quote}

\begin{quote}
Fenugreek is bitter.
\end{quote}

A feminine noun.

\begin{cousinweb}
\textbf{\mw{मेथी}} (\emph{methī}) comes from a reconstructed \emph{*metthī}, \textquotedblleft{}fenugreek\textquotedblright{}; CDIAL lists Hindi, Gujarati and Sindhi \emph{methī} (CDIAL 10313). CDIAL gives no Marwari form in this entry.
\end{cousinweb}

\begin{grammarlens}[title={What you've built}]
One more word for this chapter.
\end{grammarlens}

\subsection*{Your turn}
\begin{itemize}
\raggedright
  \item \emph{Say it:} \emph{methī}
  \item \emph{Say it:} \emph{methī}, once more
  \item \emph{Say it:} \emph{methī kaṛvī hai.} — Fenugreek is bitter.
  \item \emph{Recall:} say the Marwadi for coriander, then the Marwadi for mustard seed, then say \emph{methī} again
\end{itemize}

\begin{tcolorbox}[breakable,colback=teal!4,colframe=teal!35!black,title={Before you move on}]
What does \mw{मेथी} mean? (\textquotedblleft{}Fenugreek\textquotedblright{}.) Say it once more.
\end{tcolorbox}
\section[\texorpdfstring{\mw{हींग} (hī̃g)}{हींग (hī̃g)}]{\mw{हींग} (hī̃g) — asafoetida}

\label{lesson:MW-C185-hig}



\begin{quote}
Before the new one: say the Marwadi for mustard seed, then the Marwadi for fenugreek.
\end{quote}

\subsection*{You'll want to know: \mw{हींग}}
\textbf{\mw{हींग}} — \emph{hī̃g} — \textquotedblleft{}asafoetida\textquotedblright{}.

\begin{quote}
\mw{दाल} \mw{मांय} \mw{हींग} \mw{घालो।}
\end{quote}

\begin{quote}
\emph{dāl mā̃y hī̃g ghālo.}
\end{quote}

\begin{quote}
Put asafoetida in the dal.
\end{quote}

A feminine noun.

\begin{cousinweb}
\textbf{\mw{हींग}} (\emph{hī̃g}) is inherited from Sanskrit \emph{hiṅgu}, the asafoetida plant and the gum from its root; CDIAL lists Hindi \emph{hī̃g} (CDIAL 14079). CDIAL gives no Marwari form in this entry.
\end{cousinweb}

\begin{grammarlens}[title={What you've built}]
One more word for this chapter.
\end{grammarlens}

\subsection*{Your turn}
\begin{itemize}
\raggedright
  \item \emph{Say it:} \emph{hī̃g}
  \item \emph{Say it:} \emph{hī̃g}, once more
  \item \emph{Say it:} \emph{dāl mā̃y hī̃g ghālo.} — Put asafoetida in the dal.
  \item \emph{Recall:} say the Marwadi for mustard seed, then the Marwadi for fenugreek, then say \emph{hī̃g} again
\end{itemize}

\begin{tcolorbox}[breakable,colback=teal!4,colframe=teal!35!black,title={Before you move on}]
What does \mw{हींग} mean? (\textquotedblleft{}Asafoetida\textquotedblright{}.) Say it once more.
\end{tcolorbox}
\section[\texorpdfstring{\mw{सौंफ} (saumph)}{सौंफ (saumph)}]{\mw{सौंफ} (saumph) — fennel seed}

\label{lesson:MW-C185-saumph}



\begin{quote}
Before the new one: say the Marwadi for fenugreek, then the Marwadi for asafoetida.
\end{quote}

\subsection*{You'll want to know: \mw{सौंफ}}
\textbf{\mw{सौंफ}} — \emph{saumph} — \textquotedblleft{}fennel seed\textquotedblright{}.

\begin{quote}
\mw{सौंफ} \mw{मीठी} \mw{है।}
\end{quote}

\begin{quote}
\emph{saumph mīṭhī hai.}
\end{quote}

\begin{quote}
Fennel seed is sweet.
\end{quote}

A feminine noun.

\begin{cousinweb}
\textbf{\mw{सौंफ}} (\emph{saumph}) is inherited from Sanskrit \emph{śatapuṣpa}, \textquotedblleft{}having a hundred flowers\textquotedblright{}, a name for the plant; CDIAL lists Hindi \emph{sauph} (CDIAL 12283). CDIAL gives no Marwari form in this entry.
\end{cousinweb}

\begin{grammarlens}[title={What you've built}]
One more word for this chapter.
\end{grammarlens}

\subsection*{Your turn}
\begin{itemize}
\raggedright
  \item \emph{Say it:} \emph{saumph}
  \item \emph{Say it:} \emph{saumph}, once more
  \item \emph{Say it:} \emph{saumph mīṭhī hai.} — Fennel seed is sweet.
  \item \emph{Recall:} say the Marwadi for fenugreek, then the Marwadi for asafoetida, then say \emph{saumph} again
\end{itemize}

\begin{tcolorbox}[breakable,colback=teal!4,colframe=teal!35!black,title={Before you move on}]
What does \mw{सौंफ} mean? (\textquotedblleft{}Fennel seed\textquotedblright{}.) Say it once more.
\end{tcolorbox}
